\section{M/M/1は待ちと処理を分ける}\label{theory:design:queue}
到着がPoisson過程、処理時間が指数分布、窓口1、十分な待ち場所の定常モデルを考える。到着率\(\lambda\)、サービス率\(\mu\)は同じ時間単位へそろえる。平均処理時間がtなら\(\mu=1/t\)である。利用率\(\rho=\lambda/\mu<1\)が定常分布の条件である。出生死滅の釣合いから\(p_n=(1-\rho)\rho^n\)となり、
\[L=\sum_{n=0}^{\infty}np_n=\frac{\rho}{1-\rho},\qquad W=\frac L\lambda=\frac1{\mu-\lambda}.\]
待ちだけなら\(W_q=W-1/\mu\)、\(L_q=\lambda W_q=\rho^2/(1-\rho)\)。系内人数Lと行列人数Lqを混同しない。旧原典はLを待ち行列長と称するが、補助計算\(L/\mu\)はWqに等しく、処理時間を足す手順は正しい。名称と数式の検証を分ける。\cite{design48_queue_ref}
\section{PERTの最長経路と総余裕}\label{theory:design:pert}
先行作業が複数あれば全て終了してから後続を始める。結合点の最早時刻は\(T_E(j)=\max_i[T_E(i)+t_{ij}]\)、最遅時刻は終点から\(T_L(i)=\min_j[T_L(j)-t_{ij}]\)で求める。ダミー作業は所要0で先行関係を伝える。総余裕\(TF_{ij}=T_L(j)-T_E(i)-t_{ij}\)が0の作業はクリティカルである。
全工期は最長経路長であり、最も安い作業を無条件に縮めても工期は下がらない。複数の最長経路があれば全てを短くする。平成30年度はa3とb2が同じ22日で二つの最長枝を作るため、その作業和集合が管理対象になる。短縮前後の最長経路を再計算し、別経路への支配交代も確かめる。
\section{線形計画の実行可能性と頂点}\label{theory:design:lp}
生産量を変数にし、資源上限を不等式、利益を目的関数へ置き換える。利益係数が最大の製品だけを増やす前に、全材料と時間の制約を確認する。二変数なら制約境界の交点を求め、実行可能領域の頂点の利益を比較する。台数が整数なら連続最適解をそのまま使わず整数候補を確認する。
令和2年度は\(3x+2y\le24\)、\(x+3y\le15\)の交点\((6,3)\)が当初の最適である。利益係数変更後もこの頂点を最適にするには、隣接頂点\((0,5),(8,0)\)と比較して
\[21+6\Delta c\ge15,\qquad21+6\Delta c\ge16+8\Delta c,\qquad-1\le\Delta c\le5/2.\]
端点では最適解が複数になる。元の点が最適であることと唯一であることを区別する。全組合せ最適化に指数時間下界が証明済みと主張せず、解ける部分クラスと近似品質の条件を別に扱う。
\section{ユニバーサルデザインの利用者と場面}\label{theory:design:universal}
UDは可能な限り多くの人が特別な変更なしに使える製品・環境を目指す。1997年作業班の七原則は、公平な利用、柔軟な利用、単純で直観的な利用、認知できる情報、失敗への寛容、少ない身体的努力、接近・利用に必要な大きさと空間である。原則名を対象例へ対応させ、自動ドアの省力、速度変更の柔軟性、多形式情報、低い位置からの到達を比較する。安全のため特定利用者を一律に排除する機能を、そのままUDの例にはしない。\cite{design48_ud}
\section{分割と機能独立性}\label{theory:design:principles}
多様な要求の調和を図り、複雑系を適切な機能単位へ分割して管理する。複雑であるほど優れた設計という結論は出ない。機能間の結合を減らすことは変更影響を局所化するが、インタフェース・全体相互作用を無視できるわけではない。ばらつきは設計対象であり理想値だけを実機の保証としない。
\section{角法と図面情報}\label{theory:design:drawing}
第三角では見る側と同じ側へ図を置き、平面図は上、右側面図は右。第一角では反対側となり平面図は下、左側面図は右になる。上と下の一語が異なる候補を全体の印象で正にしない。ISOは国際標準化機構で国別のイタリア規格ではない。公開規格には第一角・第三角の節があるが有料節の全図まで読了したとはしない。\cite{design48_iso}
\section{要求性能から構造詳細と照査へ}\label{theory:design:performance}
合理的構造体を実現するため、要求性能を設定し、構造計画、構造詳細の設定、性能照査へ進む。構造詳細は部材寸法・材料等の照査に必要な情報であり、詳細設計という一般工程名と異なる。供用期間の作用と劣化を踏まえて、設定した要求が満たされるかを評価する。使用性能という一つの性能だけに全要求を限定しない。\cite{design48_mlit}
